\documentclass[12pt]{article}
\usepackage[T1]{fontenc}
\usepackage[utf8]{inputenc}
\usepackage{lmodern}
\usepackage{textcomp}
\usepackage{enumitem}
\usepackage{geometry}
\geometry{margin=1in}
\begin{document}
\begin{center}
Answer Key - Economics - Undergraduate Level Assessment 7 \
\small (Answers are ordered exactly as in the question paper.)
\end{center}

\section*{Multiple Choice}
\begin{enumerate}[label=\arabic*.]
\item \textbf{Answer:} D
\\ \textit{Options:} A) A monopoly is a company that has no control over the market.; B) A monopoly is a government-run entity that controls all aspects of a market.; C) A monopoly is a business that can set its own prices.; D) A monopoly is a firm which is the sole producer of a good.; E) A monopoly is a market structure with only two firms competing against each other.
\\ \textit{Note:} A monopoly is defined as a market structure where a single firm is the exclusive provider of a good or service, allowing it to control prices and supply without competition.
\item \textbf{Answer:} A
\\ \textit{Options:} A) maximization of the firm's value; B) minimization of costs; C) maximization of market share; D) minimization of environmental impact; E) maximization of social welfare
\\ \textit{Note:} The primary objective of management in managerial economics is the maximization of the firm's value, as this aligns with the goal of enhancing shareholder wealth and ensuring the long-term sustainability and profitability of the organization.
\item \textbf{Answer:} A
\\ \textit{Options:} A) the proportion of the good of the consumer's budget is high.; B) the number of substitute products is limited.; C) the product is unique and has no substitutes.; D) the consumer has a high income.; E) the period of time to respond to a price change is short.
\\ \textit{Note:} The price elasticity of demand is greater when the proportion of the good in the consumer's budget is high because consumers are more sensitive to price changes for goods that represent a significant portion of their expenditures, leading to larger changes in quantity demanded in response to price fluctuations.
\item \textbf{Answer:} A
\\ \textit{Options:} A) Unemployment; B) Overproduction of goods and services; C) Bankruptcy; D) Reduced consumer spending; E) Decreased interest rates
\\ \textit{Note:} At the peak of a typical business cycle, unemployment poses the greatest threat to the macroeconomy because it indicates an economy that may be overheating, leading to inflationary pressures and potential downturns if production cannot be sustained at high levels.
\item \textbf{Answer:} E
\\ \textit{Options:} A) this will increase the supply of the product.; B) this will cause the producer to decrease the supply of the product.; C) this will not affect the demand for the product now or later.; D) this will decrease the supply of the product.; E) this will increase the demand for the product.
\\ \textit{Note:} When consumers anticipate a future price increase, they are likely to buy more of the product now to avoid higher costs later, thereby increasing current demand.
\end{enumerate}

\section*{True / False}
\begin{enumerate}[label=\arabic*.]
\setcounter{enumi}{5}
\item \textbf{Answer:} False
\\ \textit{Options:} A) True; B) False
\\ \textit{Note:} The statement is false because, in the presence of a positive externality, the marginal social benefit exceeds the marginal private cost at the market quantity, indicating that the overall benefits to society are greater than the costs incurred by the individual producers.
\item \textbf{Answer:} True
\\ \textit{Options:} A) True; B) False
\\ \textit{Note:} The general to specific approach often prioritizes statistical fit over theoretical coherence, leading to models that may not adequately reflect underlying economic theories or relationships.
\item \textbf{Answer:} False
\\ \textit{Options:} A) True; B) False
\\ \textit{Note:} The GDP Deflator is considered more reliable than the Consumer Price Index because it reflects the prices of all domestically produced goods and services, while the CPI only measures the price changes of a fixed basket of consumer goods, making the GDP Deflator a broader and more comprehensive indicator of inflation in the economy.
\item \textbf{Answer:} True
\\ \textit{Options:} A) True; B) False
\\ \textit{Note:} An increase in the standard of living implies that the average income or output per person has risen, which occurs when total output grows at a faster rate than the population, leading to enhanced welfare and economic well-being.
\item \textbf{Answer:} True
\\ \textit{Options:} A) True; B) False
\\ \textit{Note:} The removal of a protective tariff on imported steel allows for greater competition and lower prices, leading to a more efficient allocation of resources as consumers can access steel at a price that reflects its true market value.
\end{enumerate}

\section*{Long Form Response}
\begin{enumerate}[label=\arabic*.]
\setcounter{enumi}{10}
\item \textbf{Answer:} In capitalist economies, price determination is primarily driven by market forces, where supply and demand dictate prices, often influenced by international regulatory bodies. In contrast, centrally planned economies rely on government directives to set prices, which can lead to inefficiencies as these prices may not reflect true market conditions. This fundamental difference affects the economic functions of prices; in capitalist systems, prices serve as signals for resource allocation and incentives for production, whereas in centrally planned economies, prices may fail to convey accurate information, leading to misallocation of resources. Consequently, the efficiency and responsiveness of capitalist economies often result in more dynamic markets, while centrally planned economies may struggle with stagnation and shortages.
\\ \textit{Note:} The answer correctly identifies that price determination in capitalist economies is based on supply and demand driven by market forces, while centrally planned economies use government directives, leading to inefficiencies and differing economic functions of prices in each system.
\item \textbf{Answer:} British socialism is characterized by its evolutionary approach, emphasizing gradual reform rather than revolutionary change. This philosophy manifests in significant government involvement in key industries, where state ownership and direction are prevalent, enabling coordinated economic planning. Additionally, the extensive government planning of land use and housing aims to address urban development and housing shortages. The establishment of a nationalized health service exemplifies the commitment to social welfare, while heavy redistributive income and estate taxes support funding for various social welfare initiatives. These characteristics reflect an economic model that prioritizes social equity and public welfare, potentially leading to a more stable and cohesive society, though they may also raise concerns about efficiency and market incentives.
\\ \textit{Note:} The answer correctly identifies British socialism's focus on gradual reform, extensive government involvement in key industries and social welfare initiatives, and highlights the economic implications of these characteristics through examples like nationalized health services and coordinated economic planning.
\end{enumerate}

\vfill
\noindent\textit{End of Answer Key}
\end{document}